% -----------------------------------------------------------
\section{Withdraw}
% -----------------------------------------------------------
	\label{WITHDRAW}
	
% % ----------------------
% \subsection{See also}
% % ----------------------

% ----------------------
\subsection{1828 American Dictionary of the English Language} % *1828
% ----------------------
	\label{WITHDRAW-1828}
	
	% -----
	\subsubsection{Transitive verb}
	% -----

		\begin{compactitem}
			\item[1] To take back; to take from.
			\item[2] To recall; to cause to retire or leave; to call back or away.
		\end{compactitem}
		
	% -----
	\subsubsection{Intransitive verb}
	% -----

		\begin{compactitem}
			\item To retire; to retreat; to quit a company or place. 
		\end{compactitem}

% ----------------------
\subsection{Oxford English Dictionary} % *OED
% ----------------------
	\label{WITHDRAW-OED}

	\begin{compactitem}
		\item[I.1.a] To take back or away (something that has been given, granted, allowed, possessed, enjoyed, or experienced).
		\item[I.1.c] To cause to decline, decrease, or disappear. 
		\item[I.2.a] To draw back, take away, remove (a thing) from its place or position.
		\item[I.2.b] To take (one's eyes, etc.) off something.
		\item[I.2.c] To remove (money) from capital, or from a bank or other place of deposit.
		\item[I.2.d] To draw (a veil, curtain, etc.) back or aside; to draw back (a bolt). Now rare.
		\item[I.3] \textit{figurative.}
		\item[I.3.a] To retract, revoke, rescind. Obsolete.
		\item[I.3.b] To remove from the scope of an inquiry, from a particular category, or the like.
		\item[I.3.c] To take back, retract (one's words, an expression). Often absol. in imperative, in parliamentary procedure, to demand the withdrawal by a member of an expression or a statement.
		\item[I.3.d] To refrain from proceeding with or prosecuting (a course of action, a proposal, etc.); to cease to support or present (a candidate, etc.).
		\item[I.4.a] To keep back or hold (one's hand); also, to withhold (a blow). Obsolete.
		\item[I.4.b] To keep back, withhold (something due, customary, or necessary); hence gen. to withhold.
		\item[I.4.c] To keep back, restrain (a person, his or her desires, etc.). Obsolete.
		\item[I.5] To draw away, deflect, divert (a person, his or her mind, etc.) from an object, pursuit, line of conduct, etc.; also, to draw, attract; to distract. Now rare.
		\item[I.6.a] To remove (a person) from a position; to cause to retire or recede; occasionally to take aside; spec. to cause (a force, troops) to retire from a position; to draw off from an engagement.
		\item[I.6.b] Law. To remove (a juror) from the panel in order to put an end to the proceedings.
		\item[II] \textit{reflexive}
		\item[II.8.a] To remove oneself from a place or position; = III.12. Now rare or archaic.
		\item[II.9] To remove oneself from a condition, sphere, society, etc.; = III.13. Now rare or archaic.
		\item[II.10] To cease, refrain; = III.14. Obsolete.
		\item[II.11] To contract. Obsolete. rare.
		\item[III] \textit{intransitive.}
		\item[III.12.a] To go away, depart, or retire from a place or position, from someone's presence, to another room or a private place, etc.; to draw back or turn aside.
		\item[III.12.b] Of combatants, troops, etc.: To retire from the field of battle or any contest, or from an advanced position.
		\item[III.12.c] Of water: To subside, ebb. Obsolete.
		\item[III.13.a] To draw away from a person; to remove oneself or retire from a society or community, from publicity, etc.; to retire from participation in or pursuit of something; to resign.
		\item[III.13.b] Of an immaterial thing, a condition, etc.: To depart, disappear.
		\item[III.14] Const. of of infinitive. To cease or refrain from, or from doing, something. Obsolete.
	\end{compactitem}
	
% % ----------------------
% \subsection{Buck's Theological Dictionary (1833)} % *BTD
% % ----------------------
	% \label{WITHDRAW-BTD}

	% \begin{compactdesc}
		% \item[PAGE\#]
	% \end{compactdesc}
	
% % ----------------------
% \subsection{Restoration Lexicon} % *RL *RESTORATION LEXICON
% % ----------------------
	% \label{WITHDRAW-OED}

	% \begin{compactitem}
		% \item[]
	% \end{compactitem}

% ----------------------
\subsection{Scriptural Usage} % *SCRIPTURES
% ----------------------
	\label{WITHDRAW-SCRIPTURES}
	
	\begin{tabular}{ll}
		Total occurrences of ``withdraw'' in the scriptures & 53 \\
	\end{tabular}

	% -----
	\subsubsection{Old Testament} % *OT
	% -----
		\label{WITHDRAW-OT}
	
		% --
		\paragraph{KJV}
		% --
			\label{WITHDRAW-OT-KJV}
	
			\begin{tabular}{ll}
				Total occurrences in the OT & 18
			\end{tabular}
			
			None of the OT uses have to do with the Spirit.
			
			% \begin{compactdesc}
				% \item[]
			% \end{compactdesc}
			
		% % --
		% \paragraph{Hebrew Bible}
		% % --
			% \label{WITHDRAW-HEBREW}
		
			% %
			% \subparagraph{\texorpdfstring{\he{sample word}}{}} % paste actual hebrew text here
			% %
				% \label{PAGA} % whatever is in step bible without periods
			
				% \begin{compactdesc}
					% \item[HALOT , PDF ] \textbf{} % Use the 1994 PDF
						% \begin{compactitem}
							% \item[1]
						% \end{compactitem}
					% \item[BDB , PDF ] \textbf{}
						% \begin{compactitem}
							% \item[1]
						% \end{compactitem}
					% \item[DCH PDF ] \textbf{} % don't include the actual page number, they repeat from volume to volume
						% \begin{compactitem}
							% \item[1]
						% \end{compactitem}
				% \end{compactdesc}
				
				% \begin{compactdesc}
					% \item[Some OT uses] \textbf{}
						% \begin{compactdesc}
							% \item[]
						% \end{compactdesc}
				% \end{compactdesc}
			
		% % --
		% \paragraph{Septuagint}
		% % --
			% \label{WITHDRAW-LXX}
		
			% %
			% \subparagraph{\texorpdfstring{\gr{sample word}}{}}
			% %
		
	% -----
	\subsubsection{New Testament} % *NT
	% -----
		\label{WITHDRAW-NT}
	
		% --
		\paragraph{KJV}
		% --
			\label{WITHDRAW-NT-KJV}		
	
			\begin{tabular}{ll}
				Total occurrences in the NT & 7
			\end{tabular}
			
			None of the NT uses have to do with the Spirit.
			
			% \begin{compactdesc}
				% \item[]
			% \end{compactdesc}
			
		% % --
		% \paragraph{Greek New Testament} % *GREEK
		% % --
			% \label{WITHDRAW-GREEK}
		
			% %
			% \subparagraph{\texorpdfstring{\gr{sample word}}{}}
			% %
				% \label{KATALLAGE}
				
				% \begin{tabular}{ll}
					% Total occurrences in the NT & 
				% \end{tabular}
				
				% \begin{compactdesc}
					% \item[Some NT uses] \textbf{} % Change to "all" if they're all listed
						% \begin{compactdesc}
							% \item[]
						% \end{compactdesc}
				% \end{compactdesc}
			
				% \begin{compactdesc}
				
					% \item[BDAG ] \phantom{.}
						% \begin{compactitem}
							% \item 
						% \end{compactitem}
						
					% \item[Brill , PDF ] \phantom{.}
						% \begin{compactitem}
							% \item 
						% \end{compactitem}
						
					% \item[CGL , PDF ] \phantom{.}
						% \begin{compactitem}
							% \item 
						% \end{compactitem}
						
					% \item[LSJ , PDF ] \phantom{.}
						% \begin{compactitem}
							% \item 
						% \end{compactitem}
						
					% \item[TDNT ] \phantom{.}
						% \begin{compactitem}
							% \item 
						% \end{compactitem}
						
					% \item[TDNTA] \phantom{.}
						% \begin{compactdesc}
							% \item 
						% \end{compactdesc}
						
					% \item[NIDNTT] \phantom{.}
						% \begin{compactdesc}
							% \item 
						% \end{compactdesc}
						
					% \item[NIDNTTE] \phantom{.}
						% \begin{compactdesc}
							% \item 
						% \end{compactdesc}
						
				% \end{compactdesc}
		
	% -----
	\subsubsection{Book of Mormon} % *BOM
	% -----
		\label{WITHDRAW-BOM}
	
		\begin{tabular}{ll}
			Total occurrences in the BOM & 19
		\end{tabular}
		
		Here are all the BOM uses that have to do with the Spirit (the others are about withdrawing people, purposes, etc.; note the use in \hyperref[ETHER-11-13]{Ether 11:13} about the prophets ``mourning'' and then withdrawing):
		
		\begin{compactdesc}
			\item[{\hyperref[MOSIAH-2-36]{Mosiah 2:36}}]
			\item[{\hyperref[ALMA-34-35]{Alma 34:35}}]
			\item[{\hyperref[HELAMAN-4-24]{Helaman 4:24}}]
			\item[{\hyperref[HELAMAN-6-35]{Helaman 6:35}}]
			\item[{\hyperref[HELAMAN-13-8]{Helaman 13:8}}]
		\end{compactdesc}
		
	% -----
	\subsubsection{Doctrine and Covenants} % *DC
	% -----
		\label{WITHDRAW-DC}
	
		\begin{tabular}{ll}
			Total occurrences in the DC & 5
		\end{tabular}
		
		\begin{compactdesc}
			\item[{\hyperref[DC19-20]{DC 19:20}}]
			\item[{\hyperref[DC88-57]{DC 88:57}}]
			\item[{\hyperref[DC121-37]{DC 121:37}}] \textbf{(2)}
		\end{compactdesc}
		
	% -----
	\subsubsection{Pearl of Great Price} % *PGP
	% -----
		\label{WITHDRAW-PGP}
	
		\begin{tabular}{ll}
			Total occurrences in the PGP & 4
		\end{tabular}
		
		\begin{compactdesc}
			\item[{\hyperref[MOSES-1-9]{Moses 1:9}}]
			\item[{\hyperref[MOSES-1-15]{Moses 1:15}}]
			\item[{\hyperref[ABRAHAM-2-12]{Abraham 2:12}}] \textbf{(2)}
		\end{compactdesc}
		
% % ----------------------
% \subsection{Key Phrases} % *KEYPHRASES *PHRASES
% % ----------------------
	% \label{WITHDRAW-PHRASES}

	% \begin{compactdesc}
	
		% \item[] \textbf{\#times} | list
			% % \begin{compactdesc}
				% % \item[Bible anchor verses] \phantom{.}
					% % \begin{compactdesc}
						% % \item[Romans 5:5] \gr{}
					% % \end{compactdesc}
				% % \item[Commentary] ---
			% % \end{compactdesc}
			
		% \item[] \textbf{\#times} | list
			% % \begin{compactdesc}
				% % \item[Bible anchor verses] \phantom{.}
					% % \begin{compactdesc}
						% % \item[Romans 5:5] \gr{}
					% % \end{compactdesc}
				% % \item[Commentary] ---
			% % \end{compactdesc}
			
	% \end{compactdesc}
	
% % ----------------------
% \subsection{Bible Dictionaries} % *DICTIONARIES
% % ----------------------
	% \label{WITHDRAW-BD}

% % ----------------------
% \subsection{Restoration Usage} % *RESTORATION
% % ----------------------
	% \label{WITHDRAW-RESTORATION}

	% % -----
	% \subsubsection{Joseph Smith} % *JS
	% % -----
		% \label{WITHDRAW-JS}
	
		% \begin{compactdesc} % Format is: 21 May 1843 HC-a, or whatever date, whatever record-keeper (if there are multiple), then the passage
			% \item[{\hyperlink{KEY}{Date/Source}}]
		% \end{compactdesc}
		
	% % -----
	% \subsubsection{Temple Ordinances}
	% % -----
		% \label{WITHDRAW-TEMPLE}
	
		% \begin{compactdesc}
			% \item[{\hyperlink{KEY}{Date/Source}}]
		% \end{compactdesc}
	
	% % -----
	% \subsubsection{Brigham Young} % *BY
	% % -----
		% \label{WITHDRAW-BY}
	
		% \begin{compactdesc} % Format is: 21 May 1843 HC-a, or whatever date, whatever record-keeper (if there are multiple), then the passage
			% \item[{\hyperlink{KEY}{Date/Source}}]
		% \end{compactdesc}
	
% ----------------------
\subsection{Commentary and Studies} % *COMMENTARY *STUDIES
% ----------------------
	\label{WITHDRAW-COMMENTARY}

	% % -----
	% \subsubsection{Key Points}
	% % -----
		% \label{WITHDRAW-KEYS}
	
		% \begin{compactitem}
			% \item
		% \end{compactitem}
		
	% -----
	\subsubsection{Withdrawing and the Spirit}
	% -----
		\label{WITHDRAW-spirit}
		
		As noted above, none of the OT and NT uses of ``withdraw'' have to do with the Spirit. But this idea appears several times in restoration scripture.
		
		A few patterns:
		
			\begin{compactdesc}
			
				\item[The Spirit withdrawing from people because of wickedness and/or refusing to repent] \phantom{.}
					\begin{compactitem}
						\item[{\hyperref[ALMA-34-35]{Alma 34:35}}]
						\item[{\hyperref[HELAMAN-4-24]{Helaman 4:24}}]
						\item[{\hyperref[HELAMAN-6-35]{Helaman 6:35}}]
						\item[{\hyperref[DC121-37]{DC 121:37}}] \phantom{.}
							\begin{compactitem}
								\item In the current text, this is a passive construction, ``when it is withdrawn,'' and so it looks like it could be the Spirit withdrawing itself or the Lord withdrawing it. But the JSP text has ``when it has withdraw,'' so I think it's really in this category.
							\end{compactitem}
						\item[{\hyperref[MOSES-1-15]{Moses 1:15}}]
					\end{compactitem}
					
				\item[The Lord withdrawing the Spirit from people because of refusing to repent] \phantom{.}
					\begin{compactitem}
						\item[{\hyperref[HELAMAN-13-8]{Helaman 13:8}}]
						\item[{\hyperref[DC19-20]{DC 19:20}}]
					\end{compactitem}
			
				\item[People withdrawing themselves from the Spirit by ignoring the commandments] \phantom{.}
					\begin{compactitem}
						\item[{\hyperref[MOSIAH-2-36]{Mosiah 2:36}}]
					\end{compactitem}
					
				\item[Other interesting uses of ``withdraw''] \phantom{.}
					\begin{compactitem}
						\item[{\hyperref[DC88-57]{DC 88:57}}]
						\item[{\hyperref[DC121-37]{DC 121:37}}] (''heavens withdraw themselves'')
						\item[{\hyperref[MOSES-1-9]{Moses 1:9}}]
						\item[{\hyperref[ABRAHAM-2-12]{Abraham 2:12}}] \textbf{(2)} \phantom{.}
							\begin{compactitem}
								\item Both of these uses in Abraham, along with the Moses 1:9 and DC 88:57, are interesting to me because they are examples of the Lord withdrawing his presence and illumination from someone even though that person has not sinned or refused to repent or otherwise done anything wrong at all. It just happened that he withdrew himself. We have no control over that, and when we lose that presence, it isn't always (or maybe even often) that we have sinned or removed ourselves from its presence. It just happens in the natural course of things.
							\end{compactitem}
					\end{compactitem}
					
			\end{compactdesc}